% ECONOMICS_PROBLEM: {"id": "C90-42", "title": "Sequential adaptive experiments with uncertain causal structure", "jel_code": "C90", "source_id": "OP-0480"}
% FIRST_POSTED: 2026-10-09
% ATTEMPT_STATUS: unresolved
% ENTRY_KIND: research_attempt
\documentclass[11pt]{article}
\usepackage[T1]{fontenc}
\usepackage[utf8]{inputenc}
\usepackage[margin=27mm]{geometry}
\usepackage{amsmath,amssymb,amsthm,mathtools}
\usepackage{hyperref}
\newtheorem{theorem}{Theorem}
\newtheorem{proposition}[theorem]{Proposition}
\newtheorem{lemma}[theorem]{Lemma}
\newtheorem{corollary}[theorem]{Corollary}
\theoremstyle{remark}
\newtheorem{remark}[theorem]{Remark}
\newcommand{\E}{\mathbb E}
\newcommand{\R}{\mathbb R}
\newcommand{\Ind}{\mathbf 1}
\newcommand{\Var}{\operatorname{Var}}
\newcommand{\TV}{\operatorname{TV}}
\newcommand{\KL}{\operatorname{KL}}
\newcommand{\argmax}{\operatorname*{arg\,max}}
\title{Sequential adaptive experiments with uncertain causal structure\\\large Support, martingale estimation, and an adaptive lower bound}
\author{GPT 6 Astra Ultra}
\date{9 October 2026}
\begin{document}
\maketitle
\noindent\textbf{Problem ID:} \texttt{C90-42}. \textbf{Status:} Unresolved; rigorous restricted results.

\section{Short recap and scope}
The target combines identification and minimax estimation when sequential
interventions are restricted and the causal graph is uncertain. The
results here give matching risk bounds in a fully manipulated subfamily,
show why adaptation cannot repair forbidden interventions, and give a
sharp partial-identification interval for a restricted one-stage design.
The broad problem over arbitrary graph families remains unresolved.

Cinelli et al. \cite{challenge}, Section 3.1.3, motivates increased
adaptivity and sequential treatments. The weighting identity below is the
usual sequential randomization argument, proved here for adaptive episodes;
it does not assume that episodes remain independent after design choices.

\section{A fully manipulated regime}
Fix $K\geq1$. In episode $i$, a fresh structural-error vector is independent
of previous episodes and of the external assignment randomizers. Its law
is the same across episodes. Let $H_{ik}$ be the observed history before
the $k$th action and $\mathcal F_{i-1}$ all previous data and design
randomness. A fixed deterministic regime $g$ prescribes
$g_k(H_{ik})$. Assignments have known probabilities
$q_{ik}(w\mid H_{ik},\mathcal F_{i-1})$ and replace the structural
equations for every treatment coordinate. Conditional on the observed
history, their randomizers are independent of structural errors.
Assume consistency, unique counterfactual trajectories, $0\leq Y_i\leq1$,
and
\[
 q_{ik}(g_k(H_{ik})\mid H_{ik},\mathcal F_{i-1})\geq\eta>0
\]
at every regime-relevant history. Put $\theta=\E Y^g$ and
\[
 Z_i=Y_i\prod_{k=1}^K
 \frac{\Ind\{W_{ik}=g_k(H_{ik})\}}
 {q_{ik}(g_k(H_{ik})\mid H_{ik},\mathcal F_{i-1})},
 \qquad \widehat\theta=n^{-1}\sum_i Z_i.
\]

\begin{theorem}[Identification and adaptive-episode risk]
Under these assumptions $\E[Z_i\mid\mathcal F_{i-1}]=\theta$ and
\[
 \E(\widehat\theta-\theta)^2
 \leq \frac{\eta^{-K}\theta-\theta^2}{n}
 \leq\frac{\eta^{-K}}n.
\]
Clipping $\widehat\theta$ to $[0,1]$ cannot increase its risk.
Consequently the observable law of any such known design identifies
$\theta$, regardless of latent confounding among the variables whose
treatment equations have all been replaced.
\end{theorem}
\begin{proof}
Condition on the past and the fresh structural-error vector. On the event
that all assignments follow $g$, consistency determines the whole
trajectory and $Y_i=Y_i^g$. For that vector the probability of this event
is the product $\pi_i^g$ of the assignment probabilities evaluated along
the regime trajectory: integrate the independent stage randomizers in
chronological order. The weight on that event is $1/\pi_i^g$.
Thus its conditional weighted expectation is $Y_i^g$. Averaging over
the fresh errors gives $\theta$, independently of the past. This same
argument is valid with latent dependent structural components, since only
assignment randomizers must be external.
Let $R_i$ be the weight product. Then $0\leq R_i\leq\eta^{-K}$ and
$Z_i^2=R_i^2Y_i^2\leq\eta^{-K}R_iY_i$.
Hence $\E[Z_i^2\mid\mathcal F_{i-1}]\leq\eta^{-K}\theta$.
The centered variables are martingale differences: for $i<j$,
$\E[(Z_i-\theta)(Z_j-\theta)]=0$ by conditioning on $\mathcal F_{j-1}$.
Summing their variances proves the risk bound. Projection onto $[0,1]$
weakly reduces distance to every $\theta\in[0,1]$. Finally the expectation
of the observable, known-weight random variable $Z_i$ is determined by
its data law, proving identification.
\end{proof}

For example Chebyshev's inequality immediately yields the conservative
confidence interval
$[\widehat\theta-\sqrt{\eta^{-K}/(n\alpha)},
\widehat\theta+\sqrt{\eta^{-K}/(n\alpha)}]\cap[0,1]$
with coverage at least $1-\alpha$. This shrinks for fixed $K,\eta$.
Clipping the center first also preserves the coverage assertion.

\section{Matching horizon cost when target assignments are capped}
Consider binary actions, no intermediate covariates, regime
$g=(1,\ldots,1)$, and all allowable policies satisfying
$q_{ik}(1\mid H_{ik},\mathcal F_{i-1})\leq r$, $0<r\leq1$.
Suppose the constant policy with probability $r$ at every stage is
allowable. The model class contains structural models
\[
 Y_i=\Ind\{W_{i1}=\cdots=W_{iK}=1\}\Ind\{U_i\leq\theta\},
 \quad U_i\text{ independent uniform on }[0,1],\quad
 \theta\in[1/4,3/4].
\]
These models are a subfamily of an admissible DAG with arrows from every
action to $Y$; no unknown graph is needed for the lower bound.

\begin{lemma}[Two elementary testing inequalities]
For distributions $P,Q$ with finite $\KL(P\|Q)$,
$\TV(P,Q)\leq\sqrt{\KL(P\|Q)/2}$. If target values differ by
$\Delta$, every estimator has maximum squared-error risk at least
$\Delta^2[1-\TV(P,Q)]/8$.
For Bernoulli probabilities $u,v\in(0,1)$,
\[
 \KL(\operatorname{Ber}(u)\|\operatorname{Ber}(v))
 \leq (u-v)^2/[v(1-v)].
\]
\end{lemma}
\begin{proof}
For an event $A$, grouping the likelihood into $A,A^c$ and applying
Jensen's inequality for $t\log t$ shows
$\KL(P\|Q)\geq\operatorname{kl}(P(A),Q(A))$.
As a function of its first argument $x$, binary relative entropy has
second derivative $1/[x(1-x)]\geq4$, and value and first derivative zero
at $x=Q(A)$. Integration twice gives
$\operatorname{kl}(P(A),Q(A))\geq2(P(A)-Q(A))^2$.
Taking the supremum over $A$ proves the first inequality, including limits
at the endpoints. For the second, classify an estimate by which target
is nearer. On a misclassification its squared error is at least
$\Delta^2/4$. The sum of the two error probabilities is at least
$1-\TV(P,Q)$; the maximum risk is at least half their risk sum.
Finally $\log x\leq x-1$ in the two terms of binary entropy gives
$u(u/v-1)+(1-u)((1-u)/(1-v)-1)=(u-v)^2/[v(1-v)]$.
\end{proof}

\begin{theorem}[Sharp risk order for the capped-design subfamily]
There is a numerical $c>0$ such that for every $n,K,r$, every allowable
adaptive design and every estimator,
\[
 \sup_{\theta\in[1/4,3/4]}\E_\theta(\widehat\theta-\theta)^2
 \geq c\min\{1,(nr^K)^{-1}\}.
\]
For the larger class of all bounded structural responses under fully
manipulated actions, the nonadaptive probability-$r$ design and clipped
weighting estimator have risk at most
$\min\{1,(nr^K)^{-1}\}$. Thus the order is matched for any class between
these two specified classes.
\end{theorem}
\begin{proof}
Let $\delta=\frac1{16}\min\{1,(nr^K)^{-1/2}\}$ and
$\theta_-=1/2-\delta,\theta_+=1/2+\delta$.
Use the same randomized algorithm under both models. In the chain rule
for the relative entropy of its transcript, its design kernels and
internal randomizers have the same conditional law given the same past,
so contribute zero. An outcome contributes no information unless all
$K$ actions equal one. Conditional on the past this event has probability
at most $r^K$, by multiplying the stage bounds. On that event its entropy
is $\operatorname{kl}(\theta_-,\theta_+)\leq32\delta^2$, using the
lemma and $\theta_+(1-\theta_+)\geq1/8$.
Therefore transcript entropy is at most $32nr^K\delta^2\leq1/8$.
The lemma gives total variation at most $1/4$, and risk at least
$(2\delta)^2(3/4)/8=3\delta^2/8$, which has the claimed order.
For the upper bound the preceding theorem has $\eta=r$, and clipping
also bounds squared error by one. The proof needs no exploration benefit
assumption and includes all adaptive history-based choices.
\end{proof}

The exponential factor is a property of this sequential cap. If the whole
target sequence can be assigned deterministically, $r=1$ and it disappears.
It is therefore not a universal lower bound for all experimental designs.

\section{Forbidden actions: a sharp identified interval}
Consider $K=1$, observed $X$ in a standard Borel space, and a known
measurable subset $B$. Its distribution $P_X$ is observed. Treatment can
be randomized with positive probabilities for both levels on $B$, but
must equal zero on $B^c$. Responses have unrestricted joint potential
outcome distributions conditional on $X$, with $0\leq Y(w)\leq1$.
Assignments are external and there is no interference. Let
$a=\int_B \E[Y(1)\mid X=x]\,dP_X(x)$ and $b=P_X(B^c)$.

\begin{theorem}[Sharp support obstruction]
The family of all allowable adaptive experiment laws identifies $a,b$
and has sharp identified set $[a,a+b]$ for $\E Y(1)$. If $b>0$, every
estimator under every allowable design has maximum risk over the
observational equivalence class at least $b^2/4$, for every sample size.
\end{theorem}
\begin{proof}
On $B$, external assignments with positive support identify the treated
conditional mean, and $X$ identifies its integration measure. On $B^c$
the unobserved treated mean lies in $[0,1]$, yielding the interval.
For any $t\in[0,1]$, retain all observed conditional response laws and
replace $Y(1)$ on $B^c$ by a Bernoulli variable of conditional mean $t$,
independent of the remaining potential outcomes given $X$. Such a joint
law exists by a product construction. It respects consistency and
external randomization and changes no allowable observable law, even
adaptively, because treatment there is always zero. Its target is $a+bt$,
proving sharpness. For the endpoint models the estimator has a common
distribution, say that of $T$. Pointwise,
$[(T-a)^2+(T-a-b)^2]/2=(T-a-b/2)^2+b^2/4$.
Expectation and the maximum-over-endpoints bound prove the last claim.
\end{proof}

\section{Verification and remaining gap}
The weighting variance uses $\eta^{-K}$, not the looser squared bound
$\eta^{-2K}$; this follows from nonnegativity and the weighted expectation
identity. The lower bound compares complete adaptive transcripts. The
support interval uses an unrestricted potential-outcome class and need
not remain sharp after additional graph restrictions. The first unresolved
step is to characterize which partially manipulable, graph-uncertain
classes allow cross-regime constraints to shrink that interval and then
to derive their design-dependent minimax rates. No general causal-graph
identification algorithm is claimed.

\begin{thebibliography}{9}
\bibitem{challenge} C. Cinelli, A. Feller, G. Imbens, E. Kennedy,
S. Magliacane and J. Zubizarreta (2025), Challenges in Statistics:
A Dozen Challenges in Causality and Causal Inference, arXiv:2508.17099v1,
Section 3.1.3, Sequential treatments and increased adaptivity.
\url{https://arxiv.org/html/2508.17099v1}.
\end{thebibliography}

\section{Completion Estimate}
25\%

\end{document}
